\slajdid{09_3-s050}
\begin{frame}[label=09_3-s050]
\alert{\textbf{PTL logička kola}}\par (pass transistor logic)
\par\medskip
Prolazna logička kola su nastala u cilju da se smanji broj potrebnih tranzistora, odnosno površina koju zauzimaju. Ovo smanjenje je naročito moguće kod selektorskih logičkih funkcija:
\par\medskip
ako je $A$ onda …, a kao nije $A$ odnosno jeste $\overline A$ onda …
\par\medskip
Osnovna ideja prolaznih logičkih kola je kontrola prenosa signala sa ulaza na izlaz putem jednog tranzistora
\par\medskip\centering\begin{tikzpicture}[x=.8cm,y=.65cm,font=\sitneoznake,>=Latex]\ctikzset{bipoles/length=.6cm}
\draw(-.35,0)--(.35,0);\draw(-.4,-.13)--(.4,-.13);\draw(-.3,0)--(-.3,.4)--(-1.2,.4)node[ocirc]{}node[left]{B};\draw(.3,0)--(.3,.4)--(2.1,.4);\draw(0,-.13)--(0,-.85)node[ocirc]{}node[below]{A};\draw(1.1,.4)node[ocirc]{}node[above]{Y};\draw(1.6,.4)to[C,l=$C_L$](1.6,-1.25)node[ground]{};\end{tikzpicture}
\par\bigskip\raggedright
podrazumevajući da je izlaz kapacitivno opterećen (često se zbog toga svrstava i u „dinamička“ logička kola, mada je princip rada ipak drugačiji).
\end{frame}
